% ============================================================
% CAPITOLO 3 - PREPROCESSING: DAI VIDEO AI PUNTI CHIAVE
% ============================================================

\section{Preprocessing: dai Video ai Punti Chiave}

Il dataset ASL Citizen è composto da video in formato MP4. Per poter addestrare un modello di Machine Learning, è necessario trasformare questi video in una rappresentazione numerica strutturata. In questa sezione viene descritto il processo di \textbf{estrazione dei landmark} (punti chiave), che converte ogni video in una sequenza di vettori numerici pronti per l'addestramento.

\subsection{L'Idea di Base}

Anziché lavorare direttamente sui pixel dei video --- un approccio computazionalmente costoso e sensibile a variazioni di sfondo, illuminazione e risoluzione --- si è scelto di rappresentare ogni frame come un insieme di \textbf{punti chiave} (\textit{landmark}) del corpo umano. Ogni punto chiave ha coordinate $(x, y, z)$ normalizzate, che descrivono la posizione di un'articolazione o un punto anatomico rilevante.

Questo approccio presenta due vantaggi fondamentali:
\begin{itemize}
    \item \textbf{Riduzione della dimensionalità:} un frame video di 640$\times$480 pixel contiene oltre 900.000 valori (RGB), mentre la rappresentazione a landmark produce soltanto 258 valori per frame.
    \item \textbf{Invarianza ambientale:} i landmark catturano esclusivamente la geometria del corpo, risultando indipendenti da sfondo, illuminazione, colore della pelle e vestiti.
\end{itemize}

\subsection{MediaPipe Holistic}

Per l'estrazione dei landmark è stata utilizzata la libreria \textbf{MediaPipe} \cite{mediapipe2020}, sviluppata da Google. In particolare, il modulo \textbf{Holistic} rileva simultaneamente tre componenti del corpo:

\begin{itemize}
    \item \textbf{Pose (corpo):} 33 landmark, ognuno con 4 valori $(x, y, z, \textit{visibility})$ = \textbf{132 valori}
    \item \textbf{Mano sinistra:} 21 landmark, ognuno con 3 valori $(x, y, z)$ = \textbf{63 valori}
    \item \textbf{Mano destra:} 21 landmark, ognuno con 3 valori $(x, y, z)$ = \textbf{63 valori}
\end{itemize}

In totale, per ogni frame si ottiene un vettore di $132 + 63 + 63 = \mathbf{258}$ \textbf{valori}. Se un componente non viene rilevato in un frame (ad esempio, una mano non visibile), il vettore corrispondente viene riempito con zeri.

\subsection{Il Processo di Estrazione}

Il Jupyter notebook di estrazione (\texttt{Extract\_Landmarks.ipynb}) elabora i video secondo i seguenti passi:

\begin{enumerate}
    \item \textbf{Caricamento del CSV:} viene letto il file \texttt{train.csv} dal dataset ASL Citizen, che contiene per ogni riga il nome del file video e il \textit{gloss} (etichetta) corrispondente.

    \item \textbf{Estrazione dal ZIP:} i video sono archiviati in un file ZIP di grandi dimensioni. Per ogni video, il worker estrae temporaneamente il file MP4.

    \item \textbf{Lettura frame per frame:} utilizzando OpenCV (\texttt{cv2.VideoCapture}), ogni frame del video viene letto e convertito da BGR a RGB (il formato richiesto da MediaPipe).

    \item \textbf{Rilevamento dei landmark:} ogni frame viene processato da MediaPipe Holistic, che restituisce le coordinate dei 75 punti chiave (33 pose + 21 mano sinistra + 21 mano destra).

    \item \textbf{Costruzione della sequenza:} i 258 valori estratti da ogni frame vengono concatenati in un array. Al termine del video, l'intera sequenza viene salvata come file \texttt{.npy} (formato binario NumPy).

    \item \textbf{Organizzazione dell'output:} i file vengono salvati in cartelle nominate con il \textit{gloss} corrispondente. Ad esempio, il video del segno ``APPLE'' viene salvato in \texttt{Dataset\_Keypoints\_Train/APPLE/}.
\end{enumerate}

La funzione principale di estrazione è la seguente:

\begin{lstlisting}[language=Python, caption={Funzione di estrazione dei landmark da un singolo frame}]
def extract_keypoints(results):
    pose = np.array([[res.x, res.y, res.z, res.visibility]
        for res in results.pose_landmarks.landmark]).flatten() \
        if results.pose_landmarks else np.zeros(132)
    lh = np.array([[res.x, res.y, res.z]
        for res in results.left_hand_landmarks.landmark]).flatten() \
        if results.left_hand_landmarks else np.zeros(63)
    rh = np.array([[res.x, res.y, res.z]
        for res in results.right_hand_landmarks.landmark]).flatten() \
        if results.right_hand_landmarks else np.zeros(63)
    return np.concatenate([pose, lh, rh])
\end{lstlisting}

\subsection{Parallelizzazione}

Data la mole di video da elaborare, il notebook utilizza il modulo \texttt{multiprocessing} di Python per distribuire il lavoro su più core della CPU. Ogni \textit{worker} elabora un video in modo indipendente, inizializzando la propria istanza di MediaPipe Holistic. Il notebook supporta anche il \textbf{resume}: se un file \texttt{.npy} esiste già, il video corrispondente viene saltato, permettendo di riprendere l'elaborazione dopo un'interruzione.

\subsection{Risultati dell'Estrazione}

L'esecuzione completa ha prodotto i seguenti risultati:

\begin{center}
\begin{tabular}{lr}
\hline
\textbf{Metrica} & \textbf{Valore} \\
\hline
Video elaborati con successo & 40.154 \\
Errori di elaborazione & 17 \\
Video senza frame validi & 31 \\
Segni distinti (cartelle) & 2.730 \\
Dimensione vettore per frame & 258 valori \\
Tempo totale di elaborazione & $\sim$7 ore 24 minuti \\
\hline
\end{tabular}
\end{center}

La struttura finale dei dati è organizzata come segue:

\begin{lstlisting}[caption={Struttura delle cartelle di output}]
Dataset_Keypoints_Train/
    APPLE/
        15890366051589533-APPLE.npy    (71 frame x 258 valori)
        654679171689897-APPLE.npy      (111 frame x 258 valori)
        ...
    HELLO/
        ...
    THANK_YOU/
        ...
    (2.730 cartelle totali)
\end{lstlisting}

Ogni file \texttt{.npy} contiene un array NumPy di forma $(\textit{n\_frame}, 258)$, dove il numero di frame varia da video a video in base alla durata del segno. Questa variabilità nella lunghezza delle sequenze sarà gestita nella fase successiva di addestramento del modello.